%%%%%%%%%%%%%%%%%%%%%%%%%%%%%%%%%%%%%%%%%%%%%%%%%%%%%%%%%%%%
%% Lecture 06
%%%%%%%%%%%%%%%%%%%%%%%%%%%%%%%%%%%%%%%%%%%%%%%%%%%%%%%%%%%%

\lecture
{06}
{Friday, 21 August 2026}
{Introduction to Probability}
{Dr. Nishant Chandgotia}

%%%%%%%%%%%%%%%%%%%%%%%%%%%%%%%%%%%%%%%%%%%%%%%%%%%%%%%%%%%%
\section{Distribution Functions and Random Variables}

Random variables are measurable maps
\[
X:
(\Omega,\mathcal{B},\mathbb{P})
\longrightarrow
(\mathbb{R},\mathcal{B}(\mathbb{R})).
\]

We define the cumulative distribution function (C.D.F.) of $X$ by

\begin{definition}{Cumulative Distribution Function}{cdf}
The \emph{cumulative distribution function} of a random variable
$X$ is the function
\[
F_X:\mathbb{R}\longrightarrow[0,1]
\]
defined by
\[
F_X(t)
=
\mathbb{P}(X\leq t).
\]
\end{definition}

\begin{remark}
\textbf{Naming.}
This is what Durrett calls the \emph{distribution function}.
\end{remark}

The C.D.F. completely determines the distribution of a real-valued
random variable.

If
\(
X\sim\mu,
\)
then the C.D.F. of $X$ is given by
\[
F_X(t)
=
\mu((-\infty,t]).
\]

Conversely, suppose that a C.D.F. $F_X$ is given. We would like to
construct the corresponding probability measure $\mu_X$ on
$\mathbb{R}$.

For $a<b$, define
\[
\mu_X((a,b])
=
F_X(b)-F_X(a).
\]

One can check that this extends to a probability measure on
$\mathbb{R}$.

%%%%%%%%%%%%%%%%%%%%%%%%%%%%%%%%%%%%%%%%%%%%%%%%%%%%%%%%%%%%
\subsection{Properties of the C.D.F.}
%%%%%%%%%%%%%%%%%%%%%%%%%%%%%%%%%%%%%%%%%%%%%%%%%%%%%%%%%%%%

Let $F_X$ be the C.D.F. of a random variable $X$.

\begin{enumerate}

    \item $F_X$ is an increasing function.

    \item
    \(
    \lim_{t\to\infty}F_X(t)=1,
    \qquad
    \lim_{t\to-\infty}F_X(t)=0.
    \)

    \item $F_X$ is right-continuous:
    \[
    \lim_{t\downarrow s}F_X(t)=F_X(s).
    \]

    Indeed,
    \[
    \begin{aligned}
    \lim_{t\downarrow s}F_X(t)
    &=
    \lim_{t\downarrow s}
    \mathbb{P}(X\leq t)=
    \mathbb{P}(X\leq s),
    \end{aligned}
    \]
    by continuity from above.

\end{enumerate}

\begin{remark}

The C.D.F. need not be left-continuous.

In fact,
\[
\mathbb{P}(X=t)>0
\]
if and only if
\[
F_X(t)-F_X(t^-)>0,
\]
where
\[
F_X(t^-)
=
\lim_{s\uparrow t}F_X(s).
\]

Moreover,
\[
F_X(t)-F_X(t^-)
=
\mathbb{P}(X=t).
\]
\end{remark}

The remaining standard properties are:

\begin{enumerate}
    \setcounter{enumi}{3}

    \item $F_X$ has at most countably many discontinuities.

    \item $F_X$ is Borel measurable.

    \item
    \(
    F_X(t^-)
    =
    \mathbb{P}(X<t).
    \)
\end{enumerate}

%%%%%%%%%%%%%%%%%%%%%%%%%%%%%%%%%%%%%%%%%%%%%%%%%%%%%%%%%%%%
\section{The Probability Measure Associated with a C.D.F.}
%%%%%%%%%%%%%%%%%%%%%%%%%%%%%%%%%%%%%%%%%%%%%%%%%%%%%%%%%%%%

We now check the construction
\[
\mu_X((a,b])
=
F_X(b)-F_X(a)
\]
and see why it defines a probability measure.

The natural collection of sets to consider first is the semiring
\[
\mathcal{S}
=
\left\{
(a,b]:a<b
\right\}.
\]

Suppose
\[
(a,b]
\subseteq
\bigcup_{i=1}^{\infty}(a_i,b_i].
\]

We want to verify the appropriate countable subadditivity/additivity
property.

For each interval,
\[
\mu_X((a_i,b_i])
=
F_X(b_i)-F_X(a_i).
\]

Hence
\[
\sum_{i=1}^{\infty}
\mu_X((a_i,b_i])
=
\sum_{i=1}^{\infty}
\left(
F_X(b_i)-F_X(a_i)
\right).
\]

The key point is the right-continuity and monotonicity of $F_X$,
which allow the interval $(a,b]$ to be approximated by finitely many
of the covering intervals.

%%%%%%%%%%%%%%%%%%%%%%%%%%%%%%%%%%%%%%%%%%%%%%%%%%%%%%%%%%%%
\subsection{Finite Additivity}
%%%%%%%%%%%%%%%%%%%%%%%%%%%%%%%%%%%%%%%%%%%%%%%%%%%%%%%%%%%%

Suppose
\[
(a,b]
=
\bigcup_{i=1}^{n}(a_i,b_i]
\]
where the intervals are disjoint and ordered so that
\[
a=a_1<b_1=a_2<b_2=a_3<\cdots<a_n<b_n=b.
\]

Then
\[
\begin{aligned}
\mu_X((a,b])
&=
F_X(b)-F_X(a)\\
&=
\sum_{i=1}^{n}
\left(
F_X(b_i)-F_X(a_i)
\right)\\
&=
\sum_{i=1}^{n}
\mu_X((a_i,b_i]).
\end{aligned}
\]

Thus $\mu_X$ is finitely additive on the semiring $\mathcal{S}$.

%%%%%%%%%%%%%%%%%%%%%%%%%%%%%%%%%%%%%%%%%%%%%%%%%%%%%%%%%%%%
\subsection{Countable Additivity}
%%%%%%%%%%%%%%%%%%%%%%%%%%%%%%%%%%%%%%%%%%%%%%%%%%%%%%%%%%%%

The same idea extends to countable disjoint unions.

Suppose
\[
(a,b]
=
\bigcup_{i=1}^{\infty}(a_i,b_i],
\]
where the intervals are pairwise disjoint.

For the finite partial unions,
\[
\bigcup_{i=1}^{n}(a_i,b_i]
\subseteq
(a,b],
\]
and finite additivity gives
\[
\sum_{i=1}^{n}
\mu_X((a_i,b_i])
\leq
\mu_X((a,b]).
\]

Passing to the limit,
\[
\sum_{i=1}^{\infty}
\mu_X((a_i,b_i])
\leq
\mu_X((a,b]).
\]

The reverse inequality follows from the covering property, together
with the continuity properties of $F_X$.

Consequently,
\[
\mu_X((a,b])
=
\sum_{i=1}^{\infty}
\mu_X((a_i,b_i]).
\]

Hence $\mu_X$ is countably additive on the semiring.

%%%%%%%%%%%%%%%%%%%%%%%%%%%%%%%%%%%%%%%%%%%%%%%%%%%%%%%%%%%%
\section{Correspondence Between C.D.F.s and Probability Measures}
%%%%%%%%%%%%%%%%%%%%%%%%%%%%%%%%%%%%%%%%%%%%%%%%%%%%%%%%%%%%

Let
\[
\mathcal{F}
=
\left\{
F:
\begin{array}{c}
F\text{ satisfies the properties}\\
\text{of a C.D.F.}
\end{array}
\right\}
\]
and let
\[
\mathcal{P}
=
\left\{
\text{probability measures on }\mathbb{R}
\right\}.
\]

The correspondence
\(
F\longmapsto\mu_F
\)
is bijective.

In other words, a probability measure on $\mathbb{R}$ determines a
unique C.D.F., and a C.D.F. determines a unique probability measure.

%%%%%%%%%%%%%%%%%%%%%%%%%%%%%%%%%%%%%%%%%%%%%%%%%%%%%%%%%%%%
\section{Constructing a Random Variable from a C.D.F.}
%%%%%%%%%%%%%%%%%%%%%%%%%%%%%%%%%%%%%%%%%%%%%%%%%%%%%%%%%%%%

\begin{question}{Constructing a Random Variable from a C.D.F.}{construct-rv-from-cdf}
Suppose we are given a function $F$ satisfying the properties of a
C.D.F. Is there a random variable $X$ such that
\[
F=F_X?
\]
\end{question}

We want to construct such a random variable explicitly.

The idea is to use the probability space
\[
([0,1],\mathcal{B}([0,1]),\lambda),
\]
where $\lambda$ denotes Lebesgue measure.

%%%%%%%%%%%%%%%%%%%%%%%%%%%%%%%%%%%%%%%%%%%%%%%%%%%%%%%%%%%%
\subsection{Case 1: $F$ is Continuous and Strictly Increasing}
%%%%%%%%%%%%%%%%%%%%%%%%%%%%%%%%%%%%%%%%%%%%%%%%%%%%%%%%%%%%

Suppose first that
\[
F:(\mathbb{R})\longrightarrow(0,1)
\]
is bijective and has no discontinuities.

\medskip
\noindent\textbf{Claim.}
We can take
\[
(\Omega,\mathcal{F},\mathbb{P})
=
([0,1],\mathcal{B}([0,1]),\lambda).
\]
\medskip
For any probability measure $\mu$, we seek an algorithmic way of
constructing
\[
X:\Omega\longrightarrow\mathbb{R}
\]
such that
\[
X\sim\mu.
\]

Define
\[
X=F^{-1}.
\]

Then, for every $t\in\mathbb{R}$,

\[
\begin{aligned}
\mathbb{P}(X\leq t)
&=
\mathbb{P}\left(F^{-1}\leq t\right)\\
&=
\mathbb{P}
\left(
\left\{
\omega:
F^{-1}(\omega)\leq t
\right\}
\right)\\
&=
\mathbb{P}
\left(
\left\{
\omega:
\omega\leq F(t)
\right\}
\right)\\
&=
\mathbb{P}\left((0,F(t)]\right)\\
&=
F(t).
\end{aligned}
\]

Therefore,
\[
F_X(t)=F(t),
\]
and hence
\[
F_X=F.
\]

%%%%%%%%%%%%%%%%%%%%%%%%%%%%%%%%%%%%%%%%%%%%%%%%%%%%%%%%%%%%
\subsection{Case 2: $F$ Has a Jump}
%%%%%%%%%%%%%%%%%%%%%%%%%%%%%%%%%%%%%%%%%%%%%%%%%%%%%%%%%%%%

\textbf{Case 2.}

Suppose $F$ has a jump at $t$.

Let
\[
a=F(t^-),
\qquad
b=F(t).
\]

Define
\[
X(\omega)
=
\begin{cases}
F^{-1}(\omega), & \omega<a,\\[4pt]
t, & a\leq\omega\leq b,\\[4pt]
F^{-1}(\omega), & \omega>b.
\end{cases}
\]

Thus the entire interval $(a,b]$ is mapped to the single point $t$.

This creates the required atom:
\[
\mathbb{P}(X=t)
=
b-a
=
F(t)-F(t^-).
\]

\begin{center}
    \textbf{Todo: Add the corresponding graph.}
\end{center}

%%%%%%%%%%%%%%%%%%%%%%%%%%%%%%%%%%%%%%%%%%%%%%%%%%%%%%%%%%%%
\subsection{Case 3: $F$ Has a Flat Portion}
%%%%%%%%%%%%%%%%%%%%%%%%%%%%%%%%%%%%%%%%%%%%%%%%%%%%%%%%%%%%

\textbf{Case 3.}

Suppose $F$ is constant on an interval $[a,b]$.

Let
\[
\theta=F(a)=F(b).
\]

Define
\[
X(\omega)
=
\begin{cases}
F^{-1}(\omega), & \omega<\theta,\\[4pt]
a, & \omega=\theta,\\[4pt]
F^{-1}(\omega), & \omega>\theta.
\end{cases}
\]

The value of $X$ at the single point $\omega=\theta$ is irrelevant
for the distribution, since
\[
\lambda(\{\theta\})=0.
\]

\begin{center}
    \textbf{Todo: Add the corresponding graph.}
\end{center}

%%%%%%%%%%%%%%%%%%%%%%%%%%%%%%%%%%%%%%%%%%%%%%%%%%%%%%%%%%%%
\subsection{Case 4: General C.D.F.}
%%%%%%%%%%%%%%%%%%%%%%%%%%%%%%%%%%%%%%%%%%%%%%%%%%%%%%%%%%%%

\textbf{Case 4.}

For a general C.D.F., define the generalized inverse by

\begin{definition}{Generalized Inverse of a C.D.F.}{generalized-inverse-cdf}
For a C.D.F. $F$, its generalized inverse is defined by
\[
F^{-1}(\omega)
=
\inf
\left\{
x\in\mathbb{R}: F(x)\geq\omega
\right\},
\qquad
\omega\in(0,1).
\]
\end{definition}
Define
\[
X(\omega)=F^{-1}(\omega).
\]

We need to show that
\[
\mathbb{P}(X\leq b)=F(b)
\]
and that $X$ is measurable.

Let
\[
A_b
=
\left\{
\omega:X(\omega)\leq b
\right\}.
\]

It is enough to show that
\[
A_b=(0,F(b)].
\]

\medskip
\noindent\textbf{Claim.}
For every $b\in\mathbb{R}$,
\[
\left\{
\omega:F^{-1}(\omega)\leq b
\right\}
=
(0,F(b)].
\]
\medskip
Consequently,
\[
\begin{aligned}
\mathbb{P}(X\leq b)
&=
\lambda(A_b)\\
&=
\lambda((0,F(b)])\\
&=
F(b).
\end{aligned}
\]

Therefore,
\[
F_X(b)=F(b)
\]
for every $b\in\mathbb{R}$.

\textbf{Homework:} Prove the claim above and complete the remaining
details of the construction.
 


%%%%%%%%%%%%%%%%%%%%%%%%%%%%%%%%%%%%%%%%%%%%%%%%%%%%%%%%%%%%

\lectureend


\section{Quiz, 24 August 2026}

Let $\mathcal{R}$ denote the respective appropriate Borel sigma algebra,
while $m$ is the Lebesgue measure on $(0,1)$.

\begin{quiz}

\begin{enumerate}

    \item Let $F:\mathbb{R}\to(0,1)$ be a function such that it is
    non-decreasing, right continuous, and
    \[
        \lim_{x\to-\infty}F(x)=0
        \qquad\text{and}\qquad
        \lim_{x\to\infty}F(x)=1.
    \]
    Prove that there is a probability measure $\mu$ on
    $(\mathbb{R},\mathcal{R})$ such that
    \[
        F(x)=\mu((-\infty,x]).
    \]

    \item Let $(\Omega,\mathcal{B},\mu)$ be a probability space and
    $X$ be a non-negative $L^1$ random variable. If $F_X$ is
    the c.d.f., prove that
    \[
        \int_{\mathbb{R}}(1-F_X(x))\,dx=\mathbb{E}(X).
    \]

    \item \textbf{This is for Bonus:} Let $\mu$ be a probability measure on
    $(\mathbb{R}^2,\mathcal{R})$. Construct a random variable
    \[
        X:((0,1),\mathcal{B},m)\to\mathbb{R}^2
    \]
    such that the distribution of $X$ is $\mu$.

\end{enumerate}

\end{quiz}
